\documentclass[12pt, a4paper]{article}

\setcounter{secnumdepth}{0}
\usepackage[margin=1in]{geometry}
\usepackage[colorlinks=true, urlcolor=blue]{hyperref}
\usepackage{titlesec}

\titleformat{\part}[block]
  {\filcenter\normalfont\Huge\bfseries}
  {}
  {0pt}
  {}

\title{Zotero's Notetaking Guide}
\author{G. Hitendra Varma}
\date{v1.0}

\begin{document}
	\maketitle
	
	\section{Introduction}
	This guide is a curation of various note-taking frameworks and methodologies designed to promote active learning and long-term knowledge retention across diverse media formats, from:
	
	\begin{itemize}
		\item Movies,
		\item Comics,
		\item Fiction and
		\item Nonfiction Books,
		\item to Academic Journals
	\end{itemize}
	
	My knowledge vault of choice is Zotero, but any note-taking app that supports colored text and highlights will work perfectly.
	
	\subsection{A Note to Readers}
	This is a highly personalized approach to consuming media. It is admittedly extensive and painstaking, often requiring multiple reviews. So even if you find this system helpful, I encourage you to adapt and modify it to fit your workflow. The core objective is to analyze a single piece of media through multiple analytical lenses.\\ Introspection is the key!
	
	\subsection{Why Zotero?}
	Comprehending the author's message and its immediate impact is only the first half of the equation; the remaining, equally crucial half is meticulous documentation and easy access, all while maintaining rigid consistency --- and there is no better archival program for this than Zotero. It's cross-platform, free, open-source, and offers generous storage on its free tier.
	
	\part{Movies}
Before we even begin to start watching a film with a critical eye, we need to do some research and document the movies' credits and specification. This is purely a archival logging, so that citations can easily be made during the actual analysis notes like \texttt{Visual Language}, \texttt{Performance}, \texttt{Audio}, \texttt{Editing}, etc.

\begin{enumerate}
	\item \textbf{Cast \& Crew}: Just a two column table with the name of the person and their role.
	
	\item \textbf{Metadata}:
	\begin{itemize}
		\item From camera models to the setup, lens series, shutter settings, exposure, aspect ratio, print format.
		
		\item Post-Production details like the final resolution, color palette, sound formats (Dolby Atmos, DTS, IMAX 6-Track), frame rate.
	\end{itemize}
\end{enumerate}

Now, the literary side of movies:

\begin{enumerate}
	\item \textbf{Characters}: Create a 
\end{enumerate}
\end{document}